% ============================================================
% Infokarten: Regierungsformen und Versorgungsmodelle
% ============================================================
\clearpage
\abschnitt{Infokarten: Regierungsformen}
Die Modelle sind bewusst vereinfacht. Sie dürfen auch Mischformen bilden, etwa einen gewählten Rat mit einer technischen Notfallleitung und einer unabhängigen Kontrolle der Sicherheitskräfte.\par\medskip

\begin{metrobox}[title={1 \ Direkte Demokratie},fonttitle=\bfseries]{cN}
Alle erwachsenen Bewohnerinnen und Bewohner stimmen direkt über wichtige Entscheidungen ab. Die Mehrheit entscheidet; grundlegende Rechte und Minderheitenschutz müssen gesondert geregelt werden.\\
\textbf{Konfliktfrage:} Ist es sinnvoll, bei einem technischen Notfall erst eine Abstimmung durchzuführen?
\end{metrobox}\par\smallskip
\begin{metrobox}[title={2 \ Repräsentative Demokratie},fonttitle=\bfseries]{cW}
Die Bevölkerung wählt einen Rat, der Entscheidungen trifft. Die fünf Gruppen könnten jeweils Vertreter entsenden, oder die Station wählt nach Einwohnerzahl.\\
\textbf{Konfliktfrage:} Soll jede Gruppe gleich viele Sitze erhalten, obwohl ihre Größe und Macht unterschiedlich sind?
\end{metrobox}\par\smallskip
\begin{metrobox}[title={3 \ Autoritäre Führung},fonttitle=\bfseries]{cS}
Eine Führungsperson oder eine kleine Führungsspitze entscheidet zentral. Sie kann schnell handeln, muss aber nicht durch die Bevölkerung legitimiert sein.\\
\textbf{Konfliktfrage:} Wie wird verhindert, dass die Führung ihre Notstandsbefugnisse dauerhaft behält?
\end{metrobox}\par\smallskip
\begin{metrobox}[title={4 \ Technokratie},fonttitle=\bfseries]{cE}
Fachleute aus Technik, Medizin und Versorgung übernehmen die Leitung. Entscheidungen orientieren sich vor allem an Sachkenntnis und Überlebensnotwendigkeiten.\\
\textbf{Konfliktfrage:} Wer kontrolliert die Fachleute, wenn ihre Entscheidungen bestimmte Gruppen benachteiligen?
\end{metrobox}\par\smallskip
\begin{metrobox}[title={5 \ Föderation (Stationskonföderation)},fonttitle=\bfseries]{cM}
Die fünf Gruppen behalten weitgehende Selbstständigkeit und schließen einen Vertrag über gemeinsame Aufgaben wie Wasserverteilung, Energieversorgung und Sicherheit.\\
\textbf{Konfliktfrage:} Was passiert, wenn eine Gruppe eine gemeinsam beschlossene Regel nicht einhält?
\end{metrobox}

\clearpage
\abschnitt{Infokarte: Versorgungsmodelle}
Zu Beginn gehören alle Marker der jeweiligen Gruppe. Im Spiel dürfen Sie beschließen, einen \textbf{gemeinsamen Vorrat} einzurichten. Die Regierungsform bestimmt, wie Sie mit den Markern umgehen: Eine Demokratie könnte über Abgaben abstimmen, eine autoritäre Führung könnte zentral beschlagnahmen, eine Technokratie könnte einen Versorgungsschlüssel festlegen.\par\medskip

{\renewcommand{\arraystretch}{1.3}\small
\begin{tabular}{|p{3.4cm}|p{6.4cm}|p{\dimexpr\textwidth-9.8cm-6\tabcolsep-4\arrayrulewidth\relax}|}
\hline
\textbf{Modell} & \textbf{Regelung} & \textbf{Mögliche Folge} \\ \hline
Freier Handel & Jede Gruppe entscheidet selbst über ihre Marker. & Hohe Eigenständigkeit, aber Gefahr ungleicher Versorgung. \\ \hline
Gemeinsamer Vorrat & Jede Gruppe zahlt einen vereinbarten Anteil ein. & Krisen können gemeinsam abgefedert werden. \\ \hline
Zentrale Verteilung & Eine Regierung entscheidet über die Verteilung. & Schnellere Entscheidungen, aber Machtmissbrauch wird möglich. \\ \hline
Verpflichtende Abgaben & Jede Gruppe muss einen festgelegten Anteil abgeben. & Die Versorgung wird planbarer, einzelne Gruppen verlieren Handlungsspielraum. \\ \hline
\end{tabular}}

\bigskip
\begin{lkbox}[title={Wichtig},fonttitle=\bfseries]
Keine Regierungsform gewinnt automatisch. Eine Regierung braucht tatsächlich vorhandene Marker: Ordnet sie Medikamente an, die nicht da sind, bleibt die Anordnung wirkungslos. Verträge sind bindend, aber nicht automatisch durchsetzbar. Wer sie bricht, zwingt die Konferenz zu einer Entscheidung.
\end{lkbox}
